%% kapa: exam=1 round=36 year=2025 period=1 form=A subject=civil_law subject_ko=민법 seq=6 number=6
%% source: Q-Net 감정평가사 기출문제 (cst00309) artlSeq=5231525 file=raw/1008/5231525/2025년 제36회 감정평가사 1차시험 1교시 문제지 원본.pdf
%% conversion: kordoc 4.24.51 -> tools/gen1.py (auto); review=proofread-sonnet; flags=none
\begin{question}{6}
법인의 등기에 관한 설명으로 옳지 않은 것은?
\begin{choices}
  \item 법인설립의 허가가 있는 때에는 3주간내에 주된 사무소소재지에서 설립등기를 해야 한다.
  \item 사원자격의 득실에 관한 사항은 사단법인의 필수등기사항이다.
  \item 청산이 종결한 때에는 청산인은 3주간내에 이를 등기해야 한다.
  \item 법인은 그 주된 사무소의 소재지에서 설립등기를 함으로써 성립한다.
  \item 이사변경등기는 제3자에 대한 대항요건이다.
\end{choices}
\end{question}
\answer{2}
